\documentclass[11pt]{article}
\usepackage[margin=1in]{geometry}
\usepackage[T1]{fontenc}
\usepackage{lmodern,amsmath,amssymb,amsthm,booktabs,array}
\usepackage[hidelinks]{hyperref}
\usepackage{microtype}
\hypersetup{pdftitle={Proof of Kimberling's replacement-iterate length and position conjectures},pdfauthor={Haobo Ma; Wenlin Zhang},pdfsubject={CC BY 4.0}}
\newtheorem{theorem}{Theorem}
\newtheorem{lemma}[theorem]{Lemma}
\newtheorem{corollary}[theorem]{Corollary}
\title{Proof of Kimberling's replacement-iterate length and position conjectures}
\author{Haobo Ma\\\small ChronoAI Pte Ltd, Singapore
\and Wenlin Zhang\\\small National University of Singapore\\\small ChronoAI Pte Ltd, Singapore}
\date{November 17, 2026}
\begin{document}
\maketitle
\begin{abstract}
We prove Clark Kimberling's length conjecture in OEIS A288476, Comments,
revision 2 (June 12, 2017), by replacing its
nonoverlapping binary rewriting rule with a three-letter substitution.
The iterate lengths satisfy $a_{n+3}=a_{n+2}+4a_{n+1}-2a_n$, starting
with $2,4,8$. The same substitution proves existence of the odd and even
limiting words and the position conjectures in A288474, A288475,
A288552 and A288553. The zero positions divided by their indices tend
to the largest root $\lambda$ of $x^3-x^2-4x+2$; the corresponding
limit for ones is $\lambda/(\lambda-1)$. All four stated position bounds hold.
For every finite binary seed we also obtain an exact length recurrence;
every seed other than the empty word and $0$ has the same pair of parity
limits, possibly interchanged.
\end{abstract}

\section{The replacement rule and the results}
Let $R$ scan a finite binary word from left to right, replacing each
original occurrence of $00$ by $0101$ and each original $1$ by $011$.
After a match it resumes after the consumed input; an unmatched $0$
is retained, and emitted output is not scanned again. This is the
\texttt{StringReplace} convention in Kimberling's program
\cite{a288473,wolfram}. Put $w_0=00$, $w_{n+1}=R(w_n)$, and
$a_n=|w_n|$, where $|w|$ denotes the length of a finite word.
Thus
\[
w_0=00,\qquad w_1=0101,\qquad w_2=00110011.
\]
The conjecture in A288476 identifies these lengths with an independently
specified recurrence \cite{a288476}. It is repeated in the entries for
the odd and even limiting words \cite{a288473,a288551}.
The earlier comment in A288473, revision 3 (June 11, 2017), refers
to A288276; Kimberling corrected that reference to A288476 in revision 5
on the same date. We use the recurrence stated explicitly in A288476.

For comparison, Dvo\v{r}\'akov\'a, Pelantov\'a and Shallit
\cite[Section 2]{dps} use a boundary-compatible block substitution for
Kimberling's different rule $00\mapsto0101$, $1\mapsto10$.
Their length sequence starts $2,4,6,10$ and has an order-four recurrence.
Here $1\mapsto011$ instead leads to a cubic recurrence and two parity
limits. We give the boundary identity explicitly, then use the same
substitution to determine densities and position bounds.

For an infinite binary word $u$, let $z_u(j)$ and $o_u(j)$ be the
one-based positions of its $j$th zero and its $j$th one, respectively,
for integers $j\geq1$.

\begin{theorem}\label{thm:main}
The lengths $a_n$ satisfy
\begin{equation}\label{eq:recurrence}
a_0=2,\quad a_1=4,\quad a_2=8,\qquad
a_{n+3}=a_{n+2}+4a_{n+1}-2a_n\quad(n\geq0),
\end{equation}
and consequently
\begin{equation}\label{eq:gf}
\sum_{n\geq0}a_n x^n=\frac{2+2x-4x^2}{1-x-4x^2+2x^3}.
\end{equation}
Both $u^{(0)}=\lim_{k\to\infty}w_{2k}$ and
$u^{(1)}=\lim_{k\to\infty}w_{2k+1}$ exist in the prefix sense.
Let $\lambda$ be the largest real root of $x^3-x^2-4x+2$, and put
$\mu=\lambda/(\lambda-1)$. For either word $u$,
\begin{equation}\label{eq:positions}
\lim_{j\to\infty}\frac{z_u(j)}j=\lambda,
\qquad
\lim_{j\to\infty}\frac{o_u(j)}j=\mu.
\end{equation}
For all $j\geq1$ and either word,
\[
-1<\lambda-\frac{z_u(j)}j<2.
\]
For $u=u^{(1)}$ the one-position bounds are
$-1<\mu-o_u(j)/j<1$; for $u=u^{(0)}$ they are
$-2<\mu-o_u(j)/j<1$.
\end{theorem}

\section{A boundary-compatible substitution}
On the alphabet $\{D,E,O\}$ define a substitution $\sigma$ and an
encoding $h$ by
\[
\begin{array}{c|ccc}
c&D&E&O\\\hline
h(c)&00&01&1\\
\sigma(c)&EE&DOO&EO
\end{array}
\]
and extend both maps to finite words by concatenation.
The binary scan $R$ itself does not preserve arbitrary concatenations:
$R(0)R(0)=00$ but $R(00)=0101$.

\begin{lemma}\label{lem:bridge}
For every letter $c\in\{D,E,O\}$ and finite binary suffix $v$,
\begin{equation}\label{eq:boundary}
R(h(c)v)=h(\sigma(c))R(v).
\end{equation}
Consequently $R(h(t))=h(\sigma(t))$ for every finite ternary word $t$,
and $w_n=h(\sigma^n(D))$ for every $n\geq0$.
\end{lemma}
\begin{proof}
For $c=D$, the initial $00$ is consumed and emits $0101=h(EE)$.
For $c=E$, the initial $0$ is unmatched and the following $1$ emits
$011$, giving $0011=h(DOO)$. For $c=O$, the initial $1$ emits
$011=h(EO)$. In every case scanning then starts at $v$.
Applying these identities to successive letters proves the word identity,
including the empty word. Since $h(D)=00$, induction on $n$ proves the last assertion.
\end{proof}

Put $A_n=|h(\sigma^n(D))|$, $B_n=|h(\sigma^n(E))|$ and
$C_n=|h(\sigma^n(O))|$. Concatenating the three substituted images gives
\begin{equation}\label{eq:states}
A_{n+1}=2B_n,\qquad B_{n+1}=A_n+2C_n,\qquad C_{n+1}=B_n+C_n.
\end{equation}
In particular $A_0=2$, $A_1=4$, $A_2=8$, and
\[
\begin{aligned}
A_{n+3}&=2A_{n+1}+4C_{n+1}
       =2A_{n+1}+4B_n+4C_n\\
       &=4A_{n+1}+4C_n
        =A_{n+2}+4A_{n+1}-2A_n.
\end{aligned}
\]
The last equality uses $A_{n+2}=2A_n+4C_n$.
Lemma~\ref{lem:bridge} identifies $A_n$ with $a_n$, proving
\eqref{eq:recurrence}. Multiplication of its formal power series by
$1-x-4x^2+2x^3$ leaves coefficients $2,2,-4$ in degrees $0,1,2$
and zero thereafter. This proves \eqref{eq:gf}, the generating function
already recorded in A288476.

\section{Parity limits and their densities}
Since
\[
\sigma^2(D)=DOODOO,\qquad \sigma^2(E)=EEEOEO,
\]
the words $\sigma^{2k}(D)$ are nested prefixes, as are
$\sigma^{2k}(E)$. Each substituted letter has at least two letters,
so their lengths tend to infinity. Their infinite limits, denoted by
$v_D$ and $v_E$, are fixed by $\sigma^2$.
The nonempty encoding preserves nested prefixes. As
$\sigma^{2k+1}(D)=\sigma^{2k}(E)\sigma^{2k}(E)$,
the binary parity limits are
$u^{(0)}=h(v_D)$ and $u^{(1)}=h(v_E)$. They are distinct:
their prefixes are $0011$ and $0101$, respectively.

The incidence matrix of $\sigma$, with columns indexed by input letters
and rows in the order $D,E,O$, is
\[
M=\begin{pmatrix}0&1&0\\2&0&1\\0&2&1\end{pmatrix},\qquad
\det(xI-M)=p(x):=x^3-x^2-4x+2.
\]
The signs of $p$ give one root in each of
\[
(-2,-3/2),\qquad (0,1/2),\qquad (7/3,12/5).
\]
These are all its roots because it has degree three. They are distinct,
and the last, $\lambda$, strictly dominates the absolute values of the others.
Positive right and left eigenvectors for $\lambda$ are
\[
r=\begin{pmatrix}1/\lambda\\1\\2/(\lambda-1)\end{pmatrix},
\qquad
\ell=\begin{pmatrix}2/\lambda\\1\\1/(\lambda-1)\end{pmatrix}.
\]
Substitution and $p(\lambda)=0$ verify $Mr=\lambda r$ and
$\ell^{T}M=\lambda\ell^{T}$. Diagonalization with the three distinct
roots therefore gives
\begin{equation}\label{eq:projection}
\lambda^{-m}M^m\longrightarrow\frac{r\ell^{T}}{\ell^{T}r}.
\end{equation}
Indeed eigenvectors for the other two roots are annihilated by
$\ell^{T}$, and their contributions tend to zero after division by
$\lambda^m$. Every column of the displayed limit is a positive multiple of $r$.

The numbers of zeros and the lengths in an encoded ternary word are
obtained by applying the row vectors $q=(2,1,0)$ and $b=(2,2,1)$ to
its letter-count column. It follows from \eqref{eq:projection} that,
uniformly over the three letters $c$,
\[
\frac{\#0\text{ in }h(\sigma^m(c))}{|h(\sigma^m(c))|}
\longrightarrow \frac{qr}{br}=\frac1\lambda.
\]
For the final equality, $qr=1+2/\lambda$ and
$br=2/\lambda+2+2/(\lambda-1)=\lambda+2$,
the latter identity being equivalent to $p(\lambda)=0$.

This block limit implies a density for every prefix of either infinite
binary word, not just for the substitution iterates. To see this, fix
$\varepsilon>0$ and take $m=2k$ so large that each of the three
encoded blocks $h(\sigma^{2k}(c))$ has zero proportion within
$\varepsilon$ of $1/\lambda$. Because $v_D$ and $v_E$ are fixed
by $\sigma^{2k}$, either infinite binary word is tiled by these blocks.
For a prefix of length $N$, the complete blocks contribute discrepancy
at most $\varepsilon N$. The one remaining partial block has length
at most $L_k=\max_c|h(\sigma^{2k}(c))|$, so its discrepancy is
at most $L_k$. Dividing by $N$, then letting $N$ tend to infinity
and $\varepsilon$ to zero, gives zero density $1/\lambda$ and
one density $(\lambda-1)/\lambda$.
Both densities are positive, so both symbols occur infinitely often.
At $N=z_u(j)$ the zero count is $j$, and at $N=o_u(j)$ the one count
is $j$. Taking reciprocals of the two positive density limits proves
\eqref{eq:positions}.

\begin{corollary}\label{cor:counts}
Let $Z_n$ and $I_n$ be the numbers of zeros and ones in $w_n$.
Each satisfies the same cubic recurrence as $a_n$, with initial triples
$(2,2,4)$ and $(0,2,4)$, respectively. Moreover,
\[
\frac{Z_n}{a_n}\longrightarrow\frac1\lambda,\qquad
\frac{I_n}{a_n}\longrightarrow\frac{\lambda-1}{\lambda},\qquad
\frac{a_{n+1}}{a_n}\longrightarrow\lambda.
\]
\end{corollary}
\begin{proof}
The three sequences are obtained by applying the rows $(2,1,0)$,
$(0,1,1)$ and $(2,2,1)$ to $M^n(1,0,0)^T$.
The characteristic identity $M^3-M^2-4M+2I=0$ gives each recurrence.
Equation~\eqref{eq:projection} gives the limits, since its column
corresponding to $D$ is a positive multiple of $r$.
\end{proof}

\section{The position inequalities}
Any finite binary input factors greedily into $00$, $01$, and $1$,
with a possible terminal $0$. The respective output blocks are
$0101$, $0011$, and $011$, with that possible terminal $0$ retained.
The complete nonterminal blocks have at least as many ones as zeros;
within each block the excess of zeros over ones is at most two,
and at a position occupied by $1$ it is at most one.
It follows that every output prefix ending at its $j$th one has
length at most $2j+1$.
Also, output begins with $0$ and contains no $111$: each block begins
with $0$ and contains at most two consecutive ones. Thus the first
zero is at position $1$ and gaps between successive zeros are at most three.
These statements pass to both parity limits, giving
\begin{equation}\label{eq:elementarybounds}
j\leq z_u(j)\leq3j-2,\qquad j\leq o_u(j)\leq2j+1.
\end{equation}
Since $7/3<\lambda<12/5$,
$12/7<\mu<7/4$. The zero-position inequalities now follow from
$1\leq z_u(j)/j\leq3$.
For the odd limit, $o_{u^{(1)}}(1)=2$, and for $j\geq2$
\eqref{eq:elementarybounds} gives $o_{u^{(1)}}(j)/j\leq5/2$.
Hence $\mu-o_{u^{(1)}}(j)/j>-1$ for every $j\geq1$.
For the even limit, $o_{u^{(0)}}(j)/j\leq3$ gives
$\mu-o_{u^{(0)}}(j)/j>-2$. In both cases the upper bound is
$\mu-o_u(j)/j\leq\mu-1<1$. This completes the proof of
Theorem~\ref{thm:main}.

\begin{corollary}\label{cor:complements}
The increasing position sequences A288474 and A288475 partition the
positive integers, as do A288552 and A288553.
\end{corollary}
\begin{proof}
At every position of either infinite binary word there is exactly one
symbol. Both symbols occur infinitely often by their positive densities.
\end{proof}

\section{The same two limits for arbitrary seeds}
The following extension concerns the same replacement rule; it is not
posed in the seven sequence entries.
\begin{corollary}\label{cor:seeds}
Every finite binary word $s$ has a unique greedy decomposition
$s=h(t)0^\delta$, where $t$ is a word over $\{D,E,O\}$ and
$\delta\in\{0,1\}$. Set $L_n=|R^n(s)|$. For every $n\geq0$,
\begin{align}
R^n(s)&=h(\sigma^n(t))0^\delta,\label{eq:seedcoding}\\
L_{n+3}&=L_{n+2}+4L_{n+1}-2L_n-2\delta,\label{eq:seedcubic}\\
L_{n+4}&=2L_{n+3}+3L_{n+2}-6L_{n+1}+2L_n.\label{eq:seedquartic}
\end{align}
If $t$ is nonempty, then the even and odd iterates of $s$ converge
in the prefix sense to $u^{(0)}$ and $u^{(1)}$, respectively, when
the first letter of $t$ is $D$ or $O$, and in the opposite order when
it is $E$. The two excluded seeds, the empty word and $0$, are fixed.
\end{corollary}
\begin{proof}
Parse $s$ from left to right as $00$, $01$, or $1$, leaving a terminal
$0$ precisely when no following symbol is available to pair with it.
The choice at every step is forced. Applying
Lemma~\ref{lem:bridge} to each letter, with suffix $0^\delta$, proves
\eqref{eq:seedcoding} by induction, since $R(0^\delta)=0^\delta$.
For the letter-count column $v$ of $t$ this gives
$L_n=(2,2,1)M^nv+\delta$. The characteristic identity for $M$
annihilates the first term; on the constant term its shift polynomial
has value $p(1)=-2$, giving \eqref{eq:seedcubic}.
Subtracting this equation at consecutive indices gives
\eqref{eq:seedquartic}.

For a first letter $D$ or $E$, the initial block
$\sigma^{2k}(D)$ or $\sigma^{2k}(E)$ has length at least $2^{2k}$
and tends to $v_D$ or $v_E$; every subsequent letter and the terminal
$0$ therefore lie beyond any fixed prefix. The odd iterates have the
other limit because $\sigma(D)=EE$ and $\sigma(E)=DOO$.
For a first letter $O$, use
$\sigma(O)=EO$ and $\sigma^2(O)=DOOEO$ to obtain the same limits
as for $D$. Encoding gives the stated binary limits.
If $t$ is empty, $s$ is the empty word or $0$, and the scanning
definition proves the final assertion directly.
\end{proof}

\section{Source coverage}

Table~\ref{tab:coverage} separates the source assertions from the
extensions above. The quoted conjectures are in the Comments of
the listed revisions; all position indices start at one.
In each position row, $m$ denotes the limit in that row.
``Proved'' refers to the written proofs above. The final section
specifies the narrower scope of the Lean declaration.
The cross-reference lists of A288552 and A288553 include A288276.
Its definition concerns rooted maps, not this replacement rule
\cite{a288276}; no assertion about that sequence follows from the
present argument.

\begin{table}[ht]
\centering\small
\begin{tabular}{@{}>{\raggedright\arraybackslash}p{.23\textwidth}>{\raggedright\arraybackslash}p{.51\textwidth}>{\raggedright\arraybackslash}p{.18\textwidth}@{}}
\toprule
Entry and locator & Statement & Coverage\\\midrule
A288476, rev. 8, Comments & ``$a(n)$ is the number of letters'' in the $n$th iterate & Proved\\
A288476, rev. 8, Formula & Recorded rational generating function & Proved\\
A288473, rev. 9, Name/Comments & Odd limiting word; repeated length conjecture & Proved\\
A288551, rev. 9, Name/Comments & Even limiting word; repeated length conjecture & Proved\\
A288474, rev. 12, Comments & ``$a(n)/n\to2.34\ldots$'' and $-1<m-a(n)/n<2$ & Proved\\
A288475, rev. 12, Comments & ``$a(n)/n\to1.74\ldots$'' and $-1<m-a(n)/n<1$ & Proved\\
A288552, rev. 8, Comments & ``$a(n)/n\to2.34\ldots$'' and $-1<m-a(n)/n<2$ & Proved\\
A288553, rev. 8, Comments & ``$a(n)/n\to1.74\ldots$'' and $-2<m-a(n)/n<1$ & Proved\\
A288474/475 and A288552/553, Names & Complementary position sequences (Corollary~\ref{cor:complements}) & Proved\\
Corollaries~\ref{cor:counts} and~\ref{cor:seeds} & Symbol counts, length growth, and arbitrary finite seeds & Not posed; proved\\
A288276, rev. 15, Name & Rooted maps of genus four with six faces & Not covered\\
\bottomrule
\end{tabular}
\caption{Source statements, with one-based position indices.}\label{tab:coverage}
\end{table}

\section{Formal scope and finite checks}
The pinned \href{https://github.com/the-omega-institute/trureturing/blob/79dfe4f731baccd42dd3c56cf6331a0ddd3cafdb/D5/S1/Words/KimberlingA288476Length.lean}{\texttt{KimberlingA288476Length.lean}}
contains the namespace
\begin{quote}\ttfamily D5.S1.Words.KimberlingA288476Length\end{quote}
Here \texttt{sourceScan} implements the original-input scanning rule,
\texttt{sourceIter} starts with $00$, and
\texttt{recurrenceSeq} is specified independently by \eqref{eq:recurrence}.
Its declaration \texttt{result} proves
\[
\bigl(|\texttt{sourceIter}(n)|:\mathbb Z\bigr)
   =\texttt{recurrenceSeq}(n)\qquad(n\in\mathbb N).
\]
The supporting declarations \texttt{scan\_code\_prefix},
\texttt{scan\_code\_word}, \texttt{native\_iterate\_code},
\texttt{coded\_step} and \texttt{coded\_recurrence} establish the
boundary identity, coding and length recurrence.
The generating function, parity limits, densities, position bounds
and Corollaries~\ref{cor:counts}--\ref{cor:seeds} have written proofs
in this note. The accompanying
\texttt{verify.py} compares direct scanning, regular-expression
replacement and ternary encoding on finite iterates, and checks the
matrix identity, rational root intervals and position inequalities
using integer and rational arithmetic. It also checks every seed
of length at most eight against \eqref{eq:seedcoding}--\eqref{eq:seedquartic}
through iterate seven and compares finite parity prefixes.
The finite checks and their
explicit ranges are recorded in \texttt{verify-output.json}.

\begingroup\small
\begin{thebibliography}{10}
\setlength{\itemsep}{2pt}
\bibitem{a288476} C. Kimberling, OEIS A288476, Comments and Formula,
revision 8 (April 6, 2020); conjecture inserted in revision 2 (June 12, 2017).
\url{https://oeis.org/A288476};
\url{https://oeis.org/history?seq=A288476}.
\bibitem{a288473} C. Kimberling, OEIS A288473, Comments and Mathematica program,
revision 9 (April 6, 2020); length comment inserted in revision 3 and
its sequence reference corrected in revision 5 (June 11, 2017).
\url{https://oeis.org/A288473}; \url{https://oeis.org/history?seq=A288473}.
\bibitem{a288474} C. Kimberling, OEIS A288474, Comments,
revision 12 (June 15, 2017). \url{https://oeis.org/A288474}.
\bibitem{a288475} C. Kimberling, OEIS A288475, Comments,
revision 12 (June 13, 2017). \url{https://oeis.org/A288475}.
\bibitem{a288551} C. Kimberling, OEIS A288551, Comments,
revision 9 (April 6, 2020). \url{https://oeis.org/A288551}.
\bibitem{a288552} C. Kimberling, OEIS A288552, Comments,
revision 8 (March 12, 2025). \url{https://oeis.org/A288552}.
\bibitem{a288553} C. Kimberling, OEIS A288553, Comments,
revision 8 (March 12, 2025). \url{https://oeis.org/A288553}.
\bibitem{a288276} G. Coserea, OEIS A288276, Name,
revision 15 (November 5, 2025). \url{https://oeis.org/A288276}.
\bibitem{dps} L. Dvo\v{r}\'akov\'a, E. Pelantov\'a and J. Shallit,
On a sequence of Kimberling and its relationship to the Tribonacci word,
\emph{Discrete Mathematics \& Theoretical Computer Science}
\textbf{28:2}, no. 19 (2026), Section 2.
\url{https://doi.org/10.46298/dmtcs.16926}; arXiv:2510.11318v5.
\bibitem{wolfram} Wolfram Research, \texttt{StringReplace}, Details,
Wolfram Language documentation. Accessed November 17, 2026.
\url{https://reference.wolfram.com/language/ref/StringReplace.html}.
\end{thebibliography}
\endgroup
\bigskip\noindent\small
\copyright\ 2026 Haobo Ma and Wenlin Zhang.
Licensed under Creative Commons Attribution 4.0 International (CC BY 4.0).
\url{https://creativecommons.org/licenses/by/4.0/}
\end{document}
